\subsection{AMR 数量的资源边际分析}

AMR 数量是仓库调度系统中最直接的\textbf{资源变量}。车辆不足时，任务会在少数 AMR 上串行排队，\textbf{时间窗延期和最大完工时间同时上升}；车辆过多时，继续扩车只能带来\textbf{有限吞吐裕度}。因此，本部分固定仓库布局、12 个任务、时间窗、电量规则和 mixed 路径代价，只改变整数资源参数 $m$，分析 $m$ 的\textbf{边际响应}。

AMR 数量 $m$ 为\textbf{离散整数资源}，局部敏感度采用\textbf{枚举重优化后的前向有限差分}。对任一随 $m$ 变化的性能指标 $G(m)$，定义
\begin{equation}
\Delta^+G(m)=G(m)-G(m+1)\qquad \text{一台 AMR 增量响应}.
\end{equation}
我们以\textbf{运筹调配问题}中最关心的\textbf{最大完工时间}为核心性能指标，并定义
\begin{equation}
\begin{aligned}
\Delta^+\Cmax(m) &= \Cmax(m)-\Cmax(m+1),\\
r_m &= \frac{\Delta^+\Cmax(m)}{\Cmax(m)},\\
\varepsilon_m &= \frac{\Delta^+\Cmax(m)/\Cmax(m)}{1/m}=m\,r_m.
\end{aligned}
\end{equation}
$\Delta^+\Cmax(m)$ 表示完工时间下降量，$r_m$ 表示相对下降比例，$\varepsilon_m$ 表示\textbf{离散弹性近似}。

为判断上述改善是否真正消除\textbf{服务风险}，结合 AMR 调运的项目具体情况构造\textbf{服务状态指标}。设 $\mathcal{J}$ 为任务集合，$\mathcal{K}_m$ 为 AMR 集合，$|\mathcal{K}_m|=m$；参数 $\theta_0$ 表示固定的布局、路径代价、时间窗、服务时间、初始位置和电量规则。给定 $m$ 后，调用主求解流程并记返回方案为 $\hat{\pi}_m$。评价器对 $\hat{\pi}_m$ 进行\textbf{时间和电量仿真}，得到任务完成时刻 $C_j(m)$、任务最晚完成时刻 $\ell_j$、路径缺失次数和电量违规次数。定义
\begin{align}
D_j(m) &= \max\{0,C_j(m)-\ell_j\}, \notag\\
\Cmax(m) &= \max_{j\in\mathcal{J}}C_j(m).
\intertext{$D_j(m)$ 表示任务延期，$\Cmax(m)$ 表示最大完工时间。据此定义服务可行性指标}
S(m) &= \sum_{j\in\mathcal{J}}D_j(m), \notag\\
L(m) &= \sum_{j\in\mathcal{J}}\indicator\{D_j(m)>0\}, \notag\\
M(m) &= N_{\mathrm{trans}}^{\mathrm{miss}}(m)+N_{\mathrm{load}}^{\mathrm{miss}}(m), \notag\\
E(m) &= N_{\mathrm{energy}}^{\mathrm{viol}}(m), \notag\\
Q(m) &= \indicator\{S(m)=0,\ L(m)=0,\ M(m)=0,\ E(m)=0\}.
\end{align}
$S(m)$ 表示总延期，$L(m)$ 表示延期任务数，$M(m)$ 表示路径不可达次数，$E(m)$ 表示电量违规次数，$Q(m)$ 表示\textbf{服务可行性}。
为辅助判断\textbf{资源是否紧张}，记 $A_{\mathrm{act}}(m)$ 为空驶、载货行驶和服务时间之和，并定义
\begin{equation}
\rho(m)=\frac{A_{\mathrm{act}}(m)}{m\,\Cmax(m)}.
\end{equation}
$\rho(m)$ 表示\textbf{平均资源利用率}。

图~\ref{fig:amr-marginal} 同时展示\textbf{资源前沿、有限差分、服务风险和车辆负载}。\textbf{4 台 AMR 是服务风险消除点}；\textbf{5 台以后 $\Cmax$ 仍下降，但边际改善明显减弱}。
\begin{center}
\includegraphics[width=0.95\linewidth]{amr_marginal}
\captionof{figure}{固定任务集下 AMR 数量的资源前沿、前向有限差分、服务风险与车辆负载}
\label{fig:amr-marginal}
\end{center}

\begin{center}
\footnotesize
\begin{minipage}[t]{0.50\linewidth}
\centering
\captionof{table}{AMR 资源前沿}
\label{tab:amr-resource-frontier}
\begin{tabular}{@{}cccccc@{}}
\toprule
$m$ & $\Cmax$ & $S$ & $L$ & $\rho$ & $Q$ \\
\midrule
2 & 84.606 & 124.593 & 4 & 97.08\% & 否 \\
3 & 59.802 & 11.802 & 1 & 91.69\% & 否 \\
4 & 40.840 & 0.000 & 0 & 95.94\% & 是 \\
5 & 34.189 & 0.000 & 0 & 92.22\% & 是 \\
6 & 29.130 & 0.000 & 0 & 90.12\% & 是 \\
\bottomrule
\end{tabular}
\end{minipage}\hfill
\begin{minipage}[t]{0.46\linewidth}
\centering
\captionof{table}{增加一台 AMR 的前向有限差分}
\label{tab:amr-finite-diff}
\begin{tabular*}{\linewidth}{@{\extracolsep{\fill}}ccccc@{}}
\toprule
区间 & $\Delta^+\Cmax$ & $r_m$ & $\varepsilon_m$ & 跨越 \\
\midrule
2--3 & 24.804 & 29.32\% & 0.586 & 否 \\
3--4 & 18.962 & 31.71\% & 0.951 & 是 \\
4--5 & 6.651 & 16.29\% & 0.651 & 否 \\
5--6 & 5.059 & 14.80\% & 0.740 & 否 \\
\bottomrule
\end{tabular*}
\end{minipage}
\end{center}

表~\ref{tab:amr-resource-frontier} 与表~\ref{tab:amr-finite-diff} 表明，\textbf{2 台和 3 台 AMR 均存在服务风险}，\textbf{$m=4$ 时总延期、延期任务数、路径缺失和电量违规同时降为 0，首次达到可行服务状态}。前向有限差分显示，\textbf{3 到 4 台是关键敏感区间}，虽然 $\Delta^+\Cmax$ 小于 2 到 3 台，但完成\textbf{服务阈值跨越}，且\textbf{离散弹性达到 0.951}。\textbf{4 台以后继续扩车只降低 $\Cmax$，不再改变服务可行性，因此 4 台 AMR 是当前算例的服务阈值，5 台可作为高裕度工况参考。}
